% GENERATED FILE. Edit canonical lessons, then run npm run generate:books.
% canonical-source-hash: fnv1a64:38af5082abc40180
% canonical-lessons: IT-C304-terremoto, IT-C304-scossa, IT-C304-macerie, IT-C304-sfollato, IT-C304-evacuare, IT-R304-ripasso-297-304

\chapter{La terra trema}
\label{ch:it-la-terra-trema}
\hlchaptermodality{304}

\begin{chapteropening}
\textbf{By the end of this chapter:} \emph{I can talk about an earthquake: the tremor, the rubble, the people who had to leave and the evacuation.}

Complete the last lesson of chapter 304: I can talk about an earthquake: the tremor, the rubble, the people who had to leave and the evacuation.
\end{chapteropening}

\section[il terremoto]{il terremoto — the earthquake}

\label{lesson:IT-C304-terremoto}



\begin{quote}
Before the new one: say the Italian for the bang, then the Italian for the collapse.
\end{quote}

\subsection*{You'll want to know: il terremoto}
\textbf{il terremoto} — \textquotedblleft{}the earthquake\textquotedblright{}.

\begin{quote}
Il terremoto è durato venti secondi.
\end{quote}

\begin{quote}
\emph{The earthquake lasted twenty seconds.}
\end{quote}

A lively child is also \emph{un terremoto}.

\begin{cousinweb}
\emph{terremoto} is inherited from Latin \emph{terrae mōtus} \textquotedblleft{}a moving of the earth\textquotedblright{} (EtymDB).
\end{cousinweb}

\begin{grammarlens}[title={What you've built}]
One more word for telling, explaining and describing at B1.
\end{grammarlens}

\subsection*{Your turn}
\begin{itemize}
\raggedright
  \item \emph{Say it:} \emph{il terremoto}
  \item \emph{Say it:} \emph{il terremoto}, once more
  \item \emph{Say it:} a sentence of your own with \emph{il terremoto}
  \item \emph{Recall:} say the Italian for the bang, then the Italian for the collapse, then say \emph{il terremoto} again
\end{itemize}

\begin{tcolorbox}[breakable,colback=teal!4,colframe=teal!35!black,title={Before you move on}]
What does il terremoto mean? (\textquotedblleft{}The earthquake\textquotedblright{}.) Say it once more.
\end{tcolorbox}
\section[la scossa]{la scossa — the tremor; the shock}

\label{lesson:IT-C304-scossa}



\begin{quote}
Before the new one: say the Italian for the collapse, then the Italian for the earthquake.
\end{quote}

\subsection*{You'll want to know: la scossa}
\textbf{la scossa} — \textquotedblleft{}the tremor; the shock\textquotedblright{}.

\begin{quote}
Stanotte abbiamo sentito una scossa.
\end{quote}

\begin{quote}
\emph{Last night we felt a tremor.}
\end{quote}

\emph{Una scossa elettrica} — an electric shock.

\begin{cousinweb}
\emph{scossa} is built on \emph{scosso}, the past participle of \emph{scuotere} \textquotedblleft{}to shake\textquotedblright{}, from Latin \emph{excussus} (EtymDB).
\end{cousinweb}

\begin{grammarlens}[title={What you've built}]
One more word for telling, explaining and describing at B1.
\end{grammarlens}

\subsection*{Your turn}
\begin{itemize}
\raggedright
  \item \emph{Say it:} \emph{la scossa}
  \item \emph{Say it:} \emph{la scossa}, once more
  \item \emph{Say it:} a sentence of your own with \emph{la scossa}
  \item \emph{Recall:} say the Italian for the collapse, then the Italian for the earthquake, then say \emph{la scossa} again
\end{itemize}

\begin{tcolorbox}[breakable,colback=teal!4,colframe=teal!35!black,title={Before you move on}]
What does la scossa mean? (\textquotedblleft{}The tremor; the shock\textquotedblright{}.) Say it once more.
\end{tcolorbox}
\section[le macerie]{le macerie — the rubble, the debris}

\label{lesson:IT-C304-macerie}



\begin{quote}
Before the new one: say the Italian for the earthquake, then the Italian for the tremor.
\end{quote}

\subsection*{You'll want to know: le macerie}
\textbf{le macerie} — \textquotedblleft{}the rubble, the debris\textquotedblright{}.

\begin{quote}
Sotto le macerie c'è ancora qualcuno.
\end{quote}

\begin{quote}
\emph{There is still someone under the rubble.}
\end{quote}

Mostly plural.

\begin{cousinweb}
\emph{macerie} is a learned borrowing from Latin \emph{māceria} \textquotedblleft{}a garden wall\textquotedblright{} (Lewis \& Short): first a dry-stone wall, later its fallen stones.
\end{cousinweb}

\begin{grammarlens}[title={What you've built}]
One more word for telling, explaining and describing at B1.
\end{grammarlens}

\subsection*{Your turn}
\begin{itemize}
\raggedright
  \item \emph{Say it:} \emph{le macerie}
  \item \emph{Say it:} \emph{le macerie}, once more
  \item \emph{Say it:} a sentence of your own with \emph{le macerie}
  \item \emph{Recall:} say the Italian for the earthquake, then the Italian for the tremor, then say \emph{le macerie} again
\end{itemize}

\begin{tcolorbox}[breakable,colback=teal!4,colframe=teal!35!black,title={Before you move on}]
What does le macerie mean? (\textquotedblleft{}The rubble, the debris\textquotedblright{}.) Say it once more.
\end{tcolorbox}
\section[lo sfollato]{lo sfollato — the evacuee, the displaced person}

\label{lesson:IT-C304-sfollato}



\begin{quote}
Before the new one: say the Italian for the tremor, then the Italian for the rubble.
\end{quote}

\subsection*{You'll want to know: lo sfollato}
\textbf{lo sfollato} — \textquotedblleft{}the evacuee, the displaced person\textquotedblright{}.

\begin{quote}
Gli sfollati dormono nella palestra della scuola.
\end{quote}

\begin{quote}
\emph{The evacuees are sleeping in the school gym.}
\end{quote}

From \emph{sfollare}, to empty of people.

\begin{cousinweb}
\emph{sfollato} is the past participle of \emph{sfollare} \textquotedblleft{}to empty of people\textquotedblright{}, built on \emph{folla} \textquotedblleft{}crowd\textquotedblright{}, which is usually traced to a Vulgar Latin \emph{fullāre} \textquotedblleft{}to tread cloth\textquotedblright{}, built on \emph{fullo} \textquotedblleft{}a fuller\textquotedblright{} (Lewis \& Short).
\end{cousinweb}

\begin{grammarlens}[title={What you've built}]
One more word for telling, explaining and describing at B1.
\end{grammarlens}

\subsection*{Your turn}
\begin{itemize}
\raggedright
  \item \emph{Say it:} \emph{lo sfollato}
  \item \emph{Say it:} \emph{lo sfollato}, once more
  \item \emph{Say it:} a sentence of your own with \emph{lo sfollato}
  \item \emph{Recall:} say the Italian for the tremor, then the Italian for the rubble, then say \emph{lo sfollato} again
\end{itemize}

\begin{tcolorbox}[breakable,colback=teal!4,colframe=teal!35!black,title={Before you move on}]
What does lo sfollato mean? (\textquotedblleft{}The evacuee, the displaced person\textquotedblright{}.) Say it once more.
\end{tcolorbox}
\section[evacuare]{evacuare — to evacuate}

\label{lesson:IT-C304-evacuare}



\begin{quote}
Before the new one: say the Italian for the rubble, then the Italian for the evacuee.
\end{quote}

\subsection*{You'll want to know: evacuare}
\textbf{evacuare} — \textquotedblleft{}to evacuate\textquotedblright{}.

\begin{quote}
Hanno evacuato tutto il paese.
\end{quote}

\begin{quote}
\emph{They evacuated the whole village.}
\end{quote}

People or buildings: \emph{evacuare la scuola}.

\begin{cousinweb}
\emph{evacuare} is a learned borrowing from Latin \emph{ēvacuāre} \textquotedblleft{}to empty out\textquotedblright{} (EtymDB), built on \emph{vacuus} \textquotedblleft{}empty\textquotedblright{}.
\end{cousinweb}

\begin{grammarlens}[title={What you've built}]
One more word for telling, explaining and describing at B1.
\end{grammarlens}

\subsection*{Your turn}
\begin{itemize}
\raggedright
  \item \emph{Say it:} \emph{evacuare}
  \item \emph{Say it:} \emph{evacuare}, once more
  \item \emph{Say it:} a sentence of your own with \emph{evacuare}
  \item \emph{Recall:} say the Italian for the rubble, then the Italian for the evacuee, then say \emph{evacuare} again
\end{itemize}

\begin{tcolorbox}[breakable,colback=teal!4,colframe=teal!35!black,title={Before you move on}]
What does evacuare mean? (\textquotedblleft{}To evacuate\textquotedblright{}.) Say it once more.
\end{tcolorbox}
\section[(dialogue)]{Ripasso dei capitoli 297–304}

\label{lesson:IT-R304-ripasso-297-304}



\begin{quote}
Say the Italian for to arrest and the arrest.
\end{quote}

\subsection*{Your turn}
Say these aloud:

\begin{itemize}
\raggedright
  \item to arrest, the arrest, the handcuffs, to capture, to flee
  \item the gang, the accomplice, to attack, the attack, to threaten
  \item the crash, to collide, to crash into something, to dent, to damage
  \item to run over, the injured person, to injure, unhurt, to come through all right
  \item emergency help, to come to the aid of, the rescuer, the stretcher, the siren
  \item the fire, to flare up, the fire extinguisher, the burn, to get burned
  \item the explosion, to explode, to burst, the bang, the collapse
  \item the earthquake, the tremor, the rubble, the evacuee, to evacuate
\end{itemize}

\begin{tcolorbox}[breakable,colback=teal!4,colframe=teal!35!black,title={Before you move on}]
To arrest? (\textbf{arrestare}.) The arrest? (\textbf{l'arresto}.) The handcuffs? (\textbf{le manette}.) To capture? (\textbf{catturare}.) To flee? (\textbf{fuggire}.) The gang? (\textbf{la banda}.) The accomplice? (\textbf{il complice}.) To attack? (\textbf{aggredire}.) The attack? (\textbf{l'aggressione}.) To threaten? (\textbf{minacciare}.) The crash? (\textbf{lo scontro}.) To collide? (\textbf{scontrarsi}.) To crash into something? (\textbf{schiantarsi}.) To dent? (\textbf{ammaccare}.) To damage? (\textbf{danneggiare}.) To run over? (\textbf{travolgere}.) The injured person? (\textbf{il ferito}.) To injure? (\textbf{ferire}.) Unhurt? (\textbf{illeso}.) To come through all right? (\textbf{cavarsela}.) Emergency help? (\textbf{il soccorso}.) To come to the aid of? (\textbf{soccorrere}.) The rescuer? (\textbf{il soccorritore}.) The stretcher? (\textbf{la barella}.) The siren? (\textbf{la sirena}.) The fire? (\textbf{l'incendio}.) To flare up? (\textbf{divampare}.) The fire extinguisher? (\textbf{l'estintore}.) The burn? (\textbf{l'ustione}.) To get burned? (\textbf{ustionarsi}.) The explosion? (\textbf{l'esplosione}.) To explode? (\textbf{esplodere}.) To burst? (\textbf{scoppiare}.) The bang? (\textbf{lo scoppio}.) The collapse? (\textbf{il crollo}.) The earthquake? (\textbf{il terremoto}.) The tremor? (\textbf{la scossa}.) The rubble? (\textbf{le macerie}.) The evacuee? (\textbf{lo sfollato}.) To evacuate? (\textbf{evacuare}.)
\end{tcolorbox}
